\subsection{Recorded Reassociation Program}
\label{app:reassociation-record}
This case connects the mathematical validation in
Proposition~\ref{prop:reassociation-program} to an emitted attribute program.
The source below has symbolic shape lists and omits all eight RHS axis
fields. The first condition preserves the ordered source batch group;
the second is the explicit representability bound used in the proposition.
The operation declaration is the generalized dot contract of
Listing~\ref{lst:dot-declaration}.

\lstinputlisting[language=TEPL,firstline=51,
caption={The exact rule used for the reassociation case study.},
label={lst:reassociation-case}]{evaluation/reassociation/reassociate.tepl}

\paragraph{Generated expressions.}
The compiler prints the outer replacement as \code{auto#2} and the inner
replacement as \code{auto#3}. Table~\ref{tab:reassociation-temporaries}
relates its temporaries to Equation~\eqref{eq:batched-program-locals}.
\code{match_positions} is a presentation of paired overlap selection;
its left/right fields are $u_G,v_G$. For the queries below, inclusion
proved in Proposition~\ref{prop:reassociation-program} makes
\code{matched_positions(query, domain)} equal to
$\pos(\mathit{domain},\mathit{query})$.

\begin{table}[!htbp]
\caption{Correspondence with the recorded compiler output.}
\label{tab:reassociation-temporaries}
\small
\begin{tabularx}{\linewidth}{@{}lY@{}}\toprule
Printed temporary & Paper expression\\\midrule
\code{m0}, \code{m2} & $i_X$, $j_Y$\\
\code{p0}, \code{p1} & $(u_I,v_I)$, $(u_J,v_J)$\\
\code{m4} & $b_Z\Vert z_J$\\
\code{m5} & $f_Y$\\
\code{m6} & $b+|f_Y|+\operatorname{range}(|f_Z|)$\\
\code{m7} & $\pos(f_Z,z_I)$\\
\code{m8} & $\gath(b+\operatorname{range}(|f_Y|),\pos(f_Y,k_Y))=s_K$\\\bottomrule
\end{tabularx}
\end{table}

In particular, the emitted outer right-contraction list is
\code{concat(gather(m6, m7), m8)}. The entrywise identity
\begin{equation}
 \gath(c+\operatorname{range}(n),q)=c+q
 \qquad\text{for an in-range selector }q
 \label{eq:shifted-range-selection}
\end{equation}
reduces it to $s_I\Vert s_K$, exactly the expression in
Table~\ref{tab:reassociation-program}. Both sides have entry $c+q_j$;
the guard bounds these integers. The remaining six direct source/selection
fields and the inner right-contraction field $z_J$ are as displayed in
that table. No inferred RHS field is treated as an input to another
attribute program.

\paragraph{Recorded execution.}
The run on 2026-10-31 used compiler SHA-256 prefix
\code{7ca6db6c4895}. Semantic checking with the algebraic policy returned
\code{Proved} and \code{additional-guard: true}; executable compilation
and serialized-program validation both exited successfully. The source,
complete compiler output, compiled JSON, command outcomes, full hashes,
checkout revision, and tracked implementation changes are retained in
\code{evaluation/reassociation/}. This is a separate case study, not an
additional success in the historical XLA totals.

The recorded normalized forms contain the same three input reads and
reduce $\vec D\Vert \vec K$ on the LHS versus $\Isel\Vert \vec K\Vert\Jsel$ on the RHS.
The proof in Appendix~\ref{app:reassociation-proof} supplies the
extent-preserving coordinate bijection and the ordered-output argument.
The emitted-program correspondence and this explicit derivation justify
the case study; the diagnostic trace is not an independently checked
certificate for all implementation reasoning.
